\section*{历史注记（第一章至第四章）}
\phantomsection
\addcontentsline{toc}{section}{第一章至第四章历史注记}

“我认为，少数选定的人能够在五年内完成此事。”——莱布尼茨（Leibniz）。{[}12 b{]}，第7卷，第187页）。

（方括号中的数字指本注记末尾的参考文献。）

对通常所谓的“数学基础”的研究，自十九世纪初以来一直坚持不懈地进行着，但若没有逻辑学系统化的并行努力——至少是逻辑学中支配数学命题之间关系的那些部分——就不可能取得成功。因此，集合论的历史和数学形式化的历史不能与“数理逻辑”的历史分开。但传统逻辑，如同现代哲学家的逻辑一样，原则上将其应用领域扩展到远比数学更广的范围。因此，读者在本注中不会找到逻辑学的历史，即使是概要形式也没有；我们尽可能地将自己限制在仅仅追溯逻辑学演变中影响数学演变的那部分。因此，我们将不谈论非经典逻辑（多值逻辑、模态逻辑）；更不用说，我们不会深入探讨从智者时代到维也纳学派时代，哲学家们就逻辑应用于外部世界对象或人类心智概念的可能性和方式所展开的争论历史。

今天已充分确立，数学在前希腊时代就已高度发达。不仅从埃及和迦勒底流传下来的最古老文献中自由使用了整数和量度大小等（极其抽象的）概念，而且巴比伦的代数学，凭借其方法的优雅和确切，不能被视为仅仅是通过经验手段解决的一堆问题。而且，尽管文本中没有任何类似于形式意义上的“证明”的东西，我们仍可以公正地认为，发现这样的解法——其普遍性从所处理的特定数值应用中显而易见——若没有最低限度的逻辑推理（也许并非完全有意识，而更像是现代代数学家在进行计算时、在写出所有细节之前所依赖的那种推理）是不可能实现的（{[}1{]}，第203页及以下）。

希腊人的独创性恰恰在于有意识地努力将数学证明按序列排列，使得从一个步骤到下一个步骤的过渡不留任何疑问，并迫使所有人一致同意。当然，希腊数学家与当今的继承者一样，在研究过程中也使用了“启发式”而非严格的论证；例如，阿基米德在他的《方法论》 {[}4 bis{]}中，提到了早期数学家“发现但未证明”的结果（*）。但是，从我们所知的最早的详细文本（可追溯至公元前五世纪中叶）起，数学文本的理想“典范”就已确定。它在欧几里得、阿基米德和阿波罗尼乌斯这些大师的著作中得到了完全实现。他们的证明观念与我们的毫无二致。

我们没有任何文本能够追溯这种“演绎方法”的最初步骤。在其首次已知出现时，它已近乎完善。我们只能认为，这是从希腊思想中持续寻求对世界的“解释”的自然发展，这种追求自公元前七世纪爱奥尼亚哲学家时代起就贯穿始终；此外，传统一致将这种方法的发展和完善归功于毕达哥拉斯学派，时间大约在公元前六世纪末至公元前五世纪中叶之间。

这种“演绎”数学，目标精确、方法严谨，将占据后世几个世纪的哲学和数学反思。一方面，我们看到“形式”逻辑的大厦以数学为模型缓慢建立，最终达到形式化语言的创造；另一方面，主要从十九世纪初开始，特别是在集合论出现之后，数学根基处的概念受到越来越密切的质疑和审视，以努力澄清其本质。

